\documentclass[11pt]{article}

% ============================================================
% Basic paper setup
% ============================================================
\usepackage[margin=1in]{geometry}
\usepackage{setspace}
\usepackage{amsmath,amssymb,amsthm}
\usepackage{mathtools}
\usepackage{bm}
\usepackage{booktabs}
\usepackage{threeparttable}
\usepackage{graphicx}
\usepackage{subcaption}
\usepackage{float}
\usepackage{placeins}
\usepackage{enumitem}
\usepackage{natbib}
\usepackage[hidelinks]{hyperref}
\usepackage[nameinlink,noabbrev]{cleveref}

% ============================================================
% Formatting
% ============================================================
\onehalfspacing
\setlength{\parindent}{1.5em}
\setlength{\parskip}{0pt}

% ============================================================
% Common math shortcuts
% ============================================================
\newcommand{\E}{\mathbb{E}}
\newcommand{\Var}{\operatorname{Var}}
\newcommand{\Cov}{\operatorname{Cov}}
\newcommand{\R}{\mathbb{R}}
\newcommand{\ind}{\mathbf{1}} % simple indicator-function notation

% ============================================================
% Theorem-style environments
% ============================================================
\newtheorem{assumption}{Assumption}
\newtheorem{proposition}{Proposition}
\newtheorem{lemma}{Lemma}

% ============================================================
% Title information
% ============================================================
\title{Paper Title}
\author{
    Your Name\\
    Department of Economics, University Name\\
    \texttt{your.email@university.edu}
}
\date{\today}

\begin{document}

\maketitle

\begin{abstract}
Write a concise abstract here. A typical economics abstract briefly states the research question, data or setting, main method, main result, and contribution.

\vspace{0.5em}
\noindent\textbf{Keywords:} keyword 1; keyword 2; keyword 3

\noindent\textbf{JEL Codes:} A10, C10
\end{abstract}

% ============================================================
\section{Introduction}\label{sec:introduction}
% ============================================================
Explain the research question, why it matters, what you do, what you find, and how the paper contributes to the literature.

A citation can look like \citet{example2020} or \citep{example2020}.

% ============================================================
\section{Institutional Background or Model}\label{sec:background}
% ============================================================
Define the setting, model, institutions, or conceptual framework.

A generic specification can be written as
\begin{equation}\label{eq:baseline}
    y_i = \alpha + \beta x_i + \gamma' z_i + \varepsilon_i,
    \qquad i=1,\ldots,n.
\end{equation}

Refer back to it as \cref{eq:baseline}.

% ============================================================
\section{Data}\label{sec:data}
% ============================================================
Describe the data source, sample construction, key variables, and summary statistics.

\begin{table}[htbp]
\centering
\caption{Summary Statistics}\label{tab:summary}
\begin{threeparttable}
\begin{tabular}{lccc}
\toprule
Variable & Mean & Std. Dev. & Observations \\
\midrule
Outcome & 1.25 & 0.42 & 1,000 \\
Treatment & 0.48 & 0.50 & 1,000 \\
Control variable & 2.10 & 1.13 & 1,000 \\
\bottomrule
\end{tabular}
\begin{tablenotes}[flushleft]
\footnotesize
\item \textit{Notes:} This table reports summary statistics for the analysis sample.
\end{tablenotes}
\end{threeparttable}
\end{table}

% ============================================================
\section{Empirical Strategy}\label{sec:strategy}
% ============================================================
Explain the identifying variation and the assumptions required for interpretation. Keep the prose primary; equations should support the argument rather than replace it.

\begin{assumption}\label{assump:example}
State an assumption clearly and in words that a reader can interpret.
\end{assumption}

% ============================================================
\section{Results}\label{sec:results}
% ============================================================
Present the main findings. Refer to tables and figures automatically, for example \cref{tab:summary}.

\begin{figure}[htbp]
\centering
% Replace the next box with \includegraphics when you have a figure file.
\fbox{\rule{0pt}{2in}\rule{0.75\linewidth}{0pt}}
\caption{Main Result}\label{fig:main}
\begin{minipage}{0.85\linewidth}
\footnotesize
\textit{Notes:} Explain what is plotted, the sample, uncertainty intervals, and any important construction details.
\end{minipage}
\end{figure}

% ============================================================
\section{Robustness and Extensions}\label{sec:robustness}
% ============================================================
Organize robustness checks by the threat or concern they address, rather than listing specifications mechanically.

% ============================================================
\section{Conclusion}\label{sec:conclusion}
% ============================================================
Summarize the main result and contribution without simply repeating the introduction.

% ============================================================
% References
% ============================================================
\bibliographystyle{apalike}
% Replace references with your .bib filename:
% \bibliography{references}

% Temporary example reference so the template compiles by itself:
\begin{thebibliography}{1}
\bibitem[Example(2020)]{example2020}
Example, A. (2020).
\newblock Example paper title.
\newblock \emph{Journal Name}, 1(1), 1--20.
\end{thebibliography}

% ============================================================
% Appendix
% ============================================================
\clearpage
\appendix
\section{Additional Results}\label{app:additional}
Put additional derivations, tables, figures, proofs, variable definitions, or robustness results here.

\end{document}
